\section{Entry from spatial initial metrics}\label{sec:spatial-entry}

The four fixed component problems on the cover have bases and bundles
\begin{equation}\label{eq:component-bundles}
 Y_j=S\times D_j,\qquad F_j=p_S^*T^{1,0}S\otimes\CC^2,
 \qquad h_{j,0}=h_S\otimes I.
\end{equation}
They have zero Einstein constant and complex parallel algebra
$I_{T^{1,0}S}\otimes\End(\CC^2)$. Their geometry is independent of
$t$ and $\theta$; the latter enters the node identifications. Thus
Theorem~\ref{thm:parallel-phase} selects, from each smooth
determinant-one component metric, a phase $I_{T^{1,0}S}\otimes P_j$.
Positivity and $(\det P_j)^2=1$ give $\det P_j=1$. For equivariant
initial profiles, uniqueness of the component flow gives
$P_{-j}=JP_jJ^*$. The selected tuple belongs to the finite-dimensional
configuration space of Section~\ref{sec:finite-flow}.

\subsection{The open class and its topology}
Put $\rho=L^{-1}$ and let $\Omega_j(2\rho)\subset D_j$ be the
complement of the two open node discs of radius $2\rho$. Use the
equivariant smooth component transfers already fixed in the
reconstruction. An actual initial metric on $\cB_{t,\theta}$ pulls
back to $h_0e^{\Psi_{\mathrm{in}}}$, where $\Psi_{\mathrm{in}}$ is a smooth
equivariant trace-free $h_0$-Hermitian endomorphism. Its restriction to
the $j$th component exterior is a section $\Psi_j$ on
$S\times\Omega_j(2\rho)$ with values in $\Herm_0(F_j)$.

Fix $0<\beta<1$ and $R_0>0$. Using the fixed component connections and
metrics in all component norms, define
\begin{equation}\label{eq:extension-norm}
 \begin{split}
 \mathcal N_t(\Psi)={}&\|\Psi\|_{C^0(X_t)}+
 \inf\left\{\sum_{j=0}^3\|\widetilde\Psi_j\|_{C^{6,\beta}(Y_j)}:
       \widetilde\Psi_j|_{S\times\Omega_j(2\rho)}=\Psi_j,\right.\\[-2pt]
 &\hspace{49mm}\left.(\widetilde\Psi_j)_{j=0}^3
                    \text{ is an equivariant smooth extension}\right\},\\
 \mathcal B_{R_0}(t,\theta)={}&
       \{h_0e^{\Psi}:\mathcal N_t(\Psi)<R_0\}.
 \end{split}
\end{equation}
Equivariance includes the base involution and the actual component
bundle lift $(i\,d\iota)\otimes J$, not just conjugation of a matrix.

\begin{lemma}\label{lem:open-class}
For every fixed smooth fibre, $\mathcal B_{R_0}(t,\theta)$ is a
nonempty open set in the $C^{6,\beta}$ topology on smooth
determinant-one metrics. It contains a parameter-independent small
ball of smooth equivariant component profiles, joined across the
nodes, and contains metrics which are not tensor products.
\end{lemma}
\begin{proof}
Each component exterior is a compact manifold with smooth boundary.
Smooth collar extension operators for its bundle sections are bounded
in the fixed $C^{6,\beta}$ norms. Averaging the operators over the
finite involution preserves equivariance, Hermitian symmetry and zero
trace. For a fixed fibre their operator bounds, together with the
triangle inequality for the infimum in \eqref{eq:extension-norm},
prove continuity of $\mathcal N_t$. The global $C^0$ term makes it
a norm, and exponentiation is smooth. No uniform extension-operator
bound as $\rho\to0$ is used. In particular the radius in
\eqref{eq:extension-norm} is not a radius for an unrestricted global
$C^0$ norm.

Take a sufficiently small fixed ball of equivariant smooth component
logarithms and join their metrics by fibrewise geodesics on the necks.
The original profiles are admissible extensions; compactness of the
positive matrix range bounds the interpolated logarithms. This yields
members of the asserted class for all small $t$. Smooth modifications
supported in the middle necks with sufficiently small $C^0$ amplitude
are allowed, independently of their derivatives.

For a non-tensor example, take a nonzero local section
$A(x)\in\Herm_0(T^{1,0}S)$ and a real smooth function $\chi(z)$
supported away from the nodes. Choose the support disjoint from its
deck image, possible because $\iota$ is free, and average the
transported logarithm $\chi(z)A(x)\otimes\sigma_z$ under the deck
action. Its mixed trace-free component remains nonzero on the original
support. Small multiples belong to the class. For positive $G,K$,
\[
 \log(G\otimes K)=(\log G)\otimes I+I\otimes(\log K),
\]
whose projection to
$\Herm_0(T^{1,0}S)\otimes\Herm_0(\CC^2)$ is zero. Uniqueness of
the Hermitian logarithm therefore excludes tensor factorization of
the example. These initial metrics can depend on all six real spatial
coordinates and need not commute with the tangent curvature.
\end{proof}

\subsection{Component evolution and the choice of phase}
Choose an equivariant extension tuple in \eqref{eq:extension-norm}
with total norm less than $2R_0$. Let $\Phi_j(s)$ be its relative
metric under the exact HYM flow on $Y_j$. Theorem~\ref{thm:parallel-phase}
and the fixed component gaps give constants depending only on $R_0$
and the fixed geometry such that
\begin{equation}\label{eq:component-phase-rate}
 \sup_{Y_j}d_4(\Phi_j(s),I\otimes P_j)^2
       \le C_{R_0}e^{-\gamma_{R_0}(s-1)},\qquad s\ge1.
\end{equation}
All these flows remain in a common compact positive range.
For every finite integer $n$, there is a common constant $C_n$
bounding their spatial covariant derivatives through order three and
their time derivatives on $0\le s\le n$. Indeed, the bounded
$C^{6,\beta}$ initial set is precompact in $C^{4,\beta/2}$ on the
fixed compact $Y_j$. The relative equation has scalar elliptic
principal part and a smooth quadratic-gradient nonlinearity on the
fixed positive range. Local parabolic continuous dependence,
smoothing and fixed-base global continuation give bounds on a
neighbourhood of each initial datum in the compact closure. A finite
cover gives $C_n$. Enlarge them to be nondecreasing and at least one.
This argument uses component derivative control, rather than drawing
derivative estimates from an arbitrary $C^0$ bound.

Two allowed extensions of the same initial metric differ only on
$S$ times discs of total area $O(\rho^2)$ and remain in a fixed
positive range. Their initial trace-distance integral is therefore
$O_{R_0}(\rho^2)$. The phase contraction
\eqref{eq:phase-contraction} implies
\begin{equation}\label{eq:phase-ambiguity}
 d_2(P_j,P'_j)^2\le C_{R_0}\rho^2.
\end{equation}
The factor two in the rank-four tensor distance and the fixed component
volume are included in the constant. Thus the phase is determined up
to $O(\rho)$, without claiming an exactly canonical extension at
finite $t$. Effective-flow nonexpansion preserves this bound at all times.

\begin{theorem}\label{thm:spatial-entry}
There is a preparation time $s_*(t)\to\infty$, with $s_*(t)/L\to0$,
independent of $\theta$ and of the initial metric in
$\mathcal B_{R_0}(t,\theta)$, with the following property.
Let $P$ be a component phase tuple selected as above. Let $H_f(s)$
solve the effective equation with $H_f(0)=P$, and let
$\widehat K(H_f(s))$ be its reconstruction. The actual HYM solution
from any metric in this open class satisfies
\begin{equation}\label{eq:entry-eta}
 \begin{gathered}
 \eta_*:=\int_{X_t}\sigma_4(h(s_*),\widehat K(P))\,dV_t
     \le C_{R_0}\left(\sqrt{\delta_L}
          +e^{-\gamma_{R_0}(s_*-1)}+L^{-2}\right)\longrightarrow0,\\
 \delta_L=\frac{1+\log L}{L}+|t|L^N,
 \end{gathered}
\end{equation}
where $N$ is a fixed finite exponent. For each $T\ge0$,
\begin{equation}\label{eq:spatial-bounded}
 \sup_{s_*\le s\le s_*+TL}\int_{X_t}
 d_4(h(s),\widehat K(H_f(s-s_*)))^2\,dV_t
       \le \eta_*+C_{R_0}Te_L.
\end{equation}
If also $\mu_0(\theta)\ge\kappa>0$, then
\begin{equation}\label{eq:spatial-alltime}
 \sup_{s\ge s_*}\int_{X_t}d_4(h(s),\widehat K(H_f(s)))^2\,dV_t
 \le C_{R_0}\eta_*+C_{R_0,\kappa}e_L(1+|\log e_L|)
                 +C_{R_0}(s_*/L)^2.
\end{equation}
All estimates before the last one are uniform through
$\theta\in\pi\ZZ$.
\end{theorem}

\begin{proof}
On the component exteriors transfer $\Phi_j(s,x,z)$ to the actual
smooth fibre. At incidence $je$ put
$A_{je}(s,x)=\Phi_j(s,x,p_{je})$. In a common node frame join the
two transported endpoints by the rank-four affine geodesic with
parameter $-L^{-1}\log|u|$. At the oriented node $j,j+1$, the
right endpoint is
\[
 (I\otimes R^*)A_{j+1,e}(I\otimes R),
\]
as follows from the actual incidence relation $w_{j+1}=Rz_j$.
Interpolate to the component metrics on $\rho<|u|<2\rho$ by
equivariant relative-logarithm cutoffs. All these constructions are
intrinsic in the positive metrics on $T^{1,0}S\otimes\CC^2$ and
respect the deck lift. They define a descended preparation metric
$K_{\mathrm{pre}}(s)$ in a fixed positive range. At $s=0$ it agrees
with the actual initial metric on the component exteriors.

We estimate its full contracted Chern residual for $0\le s\le n$.
On component regions the comparison solves the full rank-four
equation for the fixed product metric, including tangent curvature
and all spatially varying noncommuting terms. The only transfer
error is from the elliptic geometry. Smooth plumbing coefficients,
with the required two derivatives and the $C_n$ bounds, give integrated
error $C_n|t|\rho^{-N}$ for a fixed finite $N$ independent of $n$.

On a node interior, for each fixed $x\in S$ the geodesic has zero
Chern curvature in the curve direction. Its $S$-direction curvature
generally does not vanish. Covariant $S$ derivatives through order
two of the geodesic are bounded by $C_n$ on the compact endpoint
range. The tangent background has zero mean curvature, so the
$S$-contracted residual is $O(C_n)$. The time residual has the same
bound. Mixed $S$/curve curvature components can be nonzero but
contract to zero with the exact block diagonal product K\"ahler
form. Since the middle-node volume is $O(\rho^2)$, these terms
contribute $C_n\rho^2$ in integral. Thus the argument does not
assert vanishing of the full spatial curvature.

On an incidence collar the component endpoint discrepancy and the
geodesic discrepancy, with their first two curve derivatives, obey
\begin{equation}\label{eq:preparation-collar}
 C_n\bigl(\rho+L^{-1}|\log\rho|\bigr)\rho^{-k},
                    \qquad k=0,1,2.
\end{equation}
For the component term this is Taylor expansion in the curve coordinate;
for the geodesic term it is differentiation of the smooth endpoint
map at parameter $O(L^{-1}|\log\rho|)$. The relevant mixed
derivatives are included in $C_n$. Expanding the contracted Chern
expression relative to the node profile, the second derivatives of
the discrepancy and the cutoff terms integrate to
$C_n(\rho+L^{-1}|\log\rho|)$. Products with the node derivative
$O((L\rho)^{-1})$ satisfy the same bound, as do quadratic mismatch
terms for small $\rho+L^{-1}|\log\rho|$. This calculation is valid
for arbitrary rank-four matrices. The cutoff is independent of $S$,
so $S$ derivatives and time derivatives still contribute only
$C_n\rho^2$. Mixed contractions remain zero. Taking $\rho=L^{-1}$
therefore proves
\begin{equation}\label{eq:preparation-residual}
 \sup_{0\le s\le n}\int_{X_t}
       |\mathcal R(K_{\mathrm{pre}}(s))|_{\op}\,dV_t
                  \le C_n\delta_L.
\end{equation}

At time zero the exact and preparation metrics differ only in regions
of volume $O(\rho^2)$. Boundedness gives initial integrated trace
distance at most $C_{R_0}\rho^2$, regardless of the initial neck
derivatives. The forced trace comparison gives
\begin{equation}\label{eq:preparation-comparison}
 \sup_{0\le s\le n}\int_{X_t}
 d_4(h(s),K_{\mathrm{pre}}(s))^2\,dV_t
                \le C_{R_0}\rho^2+nC_n\delta_L.
\end{equation}
Choose
\begin{equation}\label{eq:preparation-time}
 s_*=\max\{n\in\mathbb N:nC_n\le\delta_L^{-1/2}\}.
\end{equation}
For small $t$ this set is nonempty and finite. Each fixed integer
eventually belongs to it, so $s_*\to\infty$. Moreover
$s_*C_{s_*}\delta_L\le\sqrt{\delta_L}$ and
$s_*/L\le\sqrt{\delta_L}$, using $C_n\ge1$ and $\delta_L\ge L^{-1}$.
This is an implicit common time; no logarithmic formula for it is claimed.

On the component exteriors at $s_*$, use
\eqref{eq:component-phase-rate}. The Green correction in
$\widehat K(P)$ has global $L^2$ size $O(L^{-1})$.
On the remaining node regions boundedness and volume $O(\rho^2)$
suffice. Therefore
\[
 \int_{X_t}d_4(K_{\mathrm{pre}}(s_*),\widehat K(P))^2\,dV_t
       \le C_{R_0}(e^{-\gamma_{R_0}(s_*-1)}+L^{-2}).
\]
Combine this with \eqref{eq:preparation-comparison} and the equivalence
of trace distance and squared affine distance on the fixed positive
range to obtain \eqref{eq:entry-eta}. Apply
Theorem~\ref{thm:actual-tracking} starting at $s_*$ to obtain
\eqref{eq:spatial-bounded}.

The effective speed is $O(L^{-1})$ for every time, so
$d_M(H_f(s-s_*),H_f(s))\le C_{R_0}s_*/L$ for $s\ge s_*$,
where $d_M$ is the mass-weighted product distance. The reconstruction
is uniformly Lipschitz into global $L^2$ on the confined tuple range.
Thus changing to the physical-time-aligned comparison costs at most
$C_{R_0}(s_*/L)^2$, with fixed squared-triangle constants. Under the
positive-rate hypothesis, Theorem~\ref{thm:all-time} applied after
entry proves \eqref{eq:spatial-alltime}.
\end{proof}

The extension ambiguity \eqref{eq:phase-ambiguity}, effective-flow nonexpansion
and the reconstruction bound show that changing the allowed component
extensions changes all subsequent comparisons by $O(L^{-1})$ in
$L^2$ distance. The theorem gives entry from a fixed-radius open class
with genuine spatial variation. It does not give trajectory tracking
during the preparation interval, uniform derivative estimates over
the whole neck, or an all-time bound as $\mu_0$ tends to zero.
